\documentclass[twocolumn,showpacs,preprintnumbers,amsmath,amssymb,prl]{revtex4-2}
% For Nature-style, use: \documentclass[12pt]{article}

\usepackage{graphicx}
\usepackage{amsmath}
\usepackage{amssymb}
\usepackage{booktabs}
\usepackage{hyperref}
\usepackage{xcolor}

\begin{document}

\title{Algebraic Structure Analysis of Fundamental Physical Constants:\\
A Computational Exploration Using Exact Rational Arithmetic}

\author{S-wind Automated Discovery System}
\affiliation{Rogue Intelligence Inc.}
\email{contact@rogueintelligence.ai}

\date{\today}

\begin{abstract}
We present a systematic computational analysis of the decimal digit sequences of fundamental physical constants using exact rational arithmetic and bounded exhaustive algebraic exploration. Our S-wind framework, employing Extended Precision Rational (EPR) numbers with zero rounding error, analyzes 10 Planck-scale constants based on CODATA 2022 recommended values \cite{codata2022}. While Planck temperature digits ($T_P = 1.416\,784(16) \times 10^{32}$ K) exhibit a quadratic fit $f(n) = \frac{7}{18}n^2 - \frac{7}{6}n + \frac{26}{15}$ that passes out-of-sample verification, we classify this as \textbf{PredictiveFit} rather than a discovery due to zero degrees of freedom in the original fit. The fine-structure constant inverse $\alpha^{-1} = 137.035\,999\,177(21)$ is classified as \textbf{MYSTERY} with a \textbf{bounded Structure-Free Certificate} under explicit search bounds ($|E| = 12$ signals, $|T| = 12$ transforms, $|H| = 19$ hypothesis classes + 130 OEIS sequences). Crucially, digit analysis yields weak, representation-dependent coincidences; algebraic analysis of Planck unit definitions yields invariant identities such as $\ell_P m_P c = \hbar$.
\end{abstract}

\pacs{02.70.-c, 06.20.Jr, 02.10.De}
\keywords{Planck constants, exact arithmetic, pattern discovery, structure-free certificate, CODATA 2022}

\maketitle

%===============================================================================
\section{Introduction}
%===============================================================================

The fundamental constants of physics---Planck's constant $h$, the speed of light $c$, the gravitational constant $G$, and the fine-structure constant $\alpha$---define the scales at which quantum mechanics, relativity, and gravity become significant \cite{planck1900,dirac1937}. A long-standing question in theoretical physics is whether these constants possess any underlying mathematical structure or are ``random'' values determined by the initial conditions of the universe \cite{barrow2002}.

Previous computational analyses have searched for patterns in the decimal expansions of mathematical constants such as $\pi$ and $e$ \cite{bailey1997}, but systematic algebraic exploration of physical constants using exact arithmetic has been limited. Floating-point arithmetic introduces rounding errors that can obscure subtle patterns or create spurious ones.

\textbf{Key insight}: For physical constants, \textit{algebraic structure analysis}---exploring relations like $\ell_P m_P c = \hbar$---is more meaningful than digit sequence analysis, which yields representation-dependent statistical coincidences at best. We demonstrate both approaches and contrast their utility.

In this work, we present a novel approach using the S-wind framework, which employs:
\begin{enumerate}
    \item \textbf{Extended Precision Rational (EPR) arithmetic}: All computations use $\frac{p}{q}$ representation with automatic GCD reduction, guaranteeing zero rounding error.
    \item \textbf{Exhaustive hypothesis testing}: Six algebraic methods (recurrence, closed-form, identity, equation solving, rational function, multiplicative structure) and 11 pattern classes are tested.
    \item \textbf{Bounded closure exhaustion}: Exploration continues until reaching closure \textit{under explicit bounds}---when no transformation can produce a new unexplored signal within the bounded search space.
    \item \textbf{Structure-Free Certificate formalization}: When no pattern is found after bounded exhaustion, the system outputs a formal certificate specifying the search bounds.
\end{enumerate}

The key innovation is the \textbf{bounded closure condition}. Let $E_S$ denote the set of explored signal signatures starting from input $S$, $T$ the transformation set, and $B$ the search bounds. We define:
\begin{equation}
\boxed{\text{CLOSURE}_B(E_S) := \forall \tau \in T, \forall S_i \in E_S: \text{sig}(\tau(S_i)) \in E_S \text{ (under bounds } B \text{)}}
\end{equation}

This guarantees that ``structure-free'' is not ``we gave up searching'' but ``we exhaustively searched $(H,T)$ under bounds $B$ and found nothing.'' The Structure-Free Certificate is then:
\begin{equation}
\text{MYSTERY}_B(S) \Leftrightarrow \text{CLOSURE}_B(E_S) \land \Big(\forall \tau \in T^{(\leq k)}, \forall f \in H: \text{DoF}(f,\tau(S))=0 \lor \text{OOS fails}\Big)
\end{equation}

%===============================================================================
\section{Methods}
%===============================================================================

\subsection{S-wind Framework}

S-wind is an exact algebraic exploration system with the following components:

\subsubsection{Extended Precision Rational (EPR)}

All numbers are represented as
\begin{equation}
    r = \frac{p}{q}, \quad p, q \in \mathbb{Z}_{128}, \quad \gcd(p, q) = 1
\end{equation}
where $\mathbb{Z}_{128}$ denotes 128-bit signed integers. This guarantees:
\begin{itemize}
    \item Zero rounding error for $+, -, \times, \div$
    \item Exact comparison operations
    \item Deterministic reproducibility
\end{itemize}

\subsubsection{Hypothesis Space $H$}

The system tests the following algebraic structures:
\begin{align}
    H = \{ &H_{\text{recurrence}}: s_{n+k} = \sum_{i=0}^{k-1} a_i s_{n+i}, \\
           &H_{\text{polynomial}}: f(n) = \sum_{i=0}^{d} a_i n^i, \\
           &H_{\text{rational}}: f(n) = \frac{P(n)}{Q(n)}, \\
           &H_{\text{multiplicative}}: f(mn) = f(m)f(n), \\
           &H_{\text{exponential}}: f(n) = a \cdot r^n, \\
           &H_{\text{periodic}}: f(n+T) = f(n) \}
\end{align}

Additionally, S-wind queries a fixed OEIS subset: \textbf{130 sequences} shipped with the system (\texttt{s-wind v1.1.1, knowledge\_base.rs}), selected by exact prefix match (first 5 terms). This subset is pre-registered and version-controlled (commit-locked at release), not selected post-hoc, eliminating data selection bias.

\subsubsection{Transformation Set $T$}

Signal transformations applied iteratively with \textbf{bounded closure}---the system continues until closure under bounds $B$ (no transform can produce a new signal within the bounded search space):
\begin{align}
    T = \{ &\tau_{\text{id}}: S \mapsto S, \\
           &\tau_{\Delta}: S \mapsto \{s_{n+1} - s_n\}, \\
           &\tau_{\Delta^2}: S \mapsto \text{second difference}, \\
           &\tau_{\Delta^3}: S \mapsto \text{third difference}, \\
           &\tau_{\text{ratio}}: S \mapsto \{s_{n+1}/s_n\}, \\
           &\tau_{\Sigma}: S \mapsto \{\sum_{k=1}^{n} s_k\}, \\
           &\tau_{|\cdot|}: S \mapsto \{|s_n|\}, \\
           &\tau_{\text{sign}}: S \mapsto \{\text{sgn}(s_n)\}, \\
           &\tau_{(\cdot)^2}: S \mapsto \{s_n^2\}, \\
           &\tau_{n \cdot f}: S \mapsto \{n \cdot s_n\}, \\
           &\tau_{\sqrt{|\cdot|}}: S \mapsto \{\sqrt{|s_n|}\}, \\
           &\tau_{\text{MA}}: S \mapsto \text{moving average}, \\
           &\tau_{\text{prod}}: S \mapsto \{s_n \cdot s_{n+1}\} \}
\end{align}

\textbf{Bounded Closure Exhaustion}: The system explores all compositions of transforms until reaching a \textit{closure}---the set of reachable signals becomes fixed. This ensures no algebraic structure is missed within $(H, T)$.

\textbf{Explicit Search Bounds $B$} (for Structure-Free Certificate validity):
\begin{itemize}
    \item Max polynomial degree: $d_{\max} = \min(|S|-1, 10)$
    \item Max recurrence order: $k_{\max} = 3$
    \item Max transform depth: $\tau_{\max} = |S|$ (full difference chain)
    \item Moving-average windows: $W = \{3\}$ (only MA3 is applied)
    \item EPR complexity cap: denominators bounded by $q_{\max} = 2^{127}$ (128-bit signed integers)
    \item Termination: Until closure reached under above bounds
\end{itemize}

These bounds guarantee that the explored signal set $E_S$ is finite and closure is always reachable in finite time.

\subsubsection{Discovery Assessment}

For each detected pattern $f$ with parameter vector $\Theta_f$, we compute:
\begin{equation}
    \text{DoF}(f, S) = |S| - \dim(\Theta_f)
\end{equation}
where $|S|$ is the sample size and $\dim(\Theta_f)$ is the number of free parameters.

\textbf{Critical insight}: By the polynomial interpolation theorem, $n$ points can \textit{always} be fitted exactly by a polynomial of degree $n-1$. Therefore:
\begin{itemize}
    \item $n$ points with $n$ parameters $\Rightarrow$ DoF $= 0$ (guaranteed fit)
    \item DoF $= 0$ carries \textbf{zero information gain}
\end{itemize}

Classification rules:
\begin{itemize}
    \item \textbf{EmergentLaw}: DoF $\geq 2$ AND OOS passes (genuine discovery)
    \item \textbf{DescriptiveModel}: DoF $> 0$ but OOS fails
    \item \textbf{PredictiveFit}: DoF $= 0$ but OOS passes (may be coincidence)
    \item \textbf{InterpolantFit}: DoF $= 0$ and OOS fails (pure interpolation)
    \item \textbf{MYSTERY}: All patterns exhausted, no structure found
\end{itemize}

\subsection{Data Source: CODATA 2022}

All physical constant values are taken from \textbf{CODATA 2022} (NIST) recommended values \cite{codata2022}:

\begin{table}[h]
\centering
\caption{Physical constants analyzed (CODATA 2022)}
\label{tab:constants}
\begin{tabular}{@{}lll@{}}
\toprule
Constant & Symbol & Value \\
\midrule
Planck constant & $h$ & $6.62607015 \times 10^{-34}$ J·s (exact) \\
Planck length & $\ell_P$ & $1.616\,255(18) \times 10^{-35}$ m \\
Planck time & $t_P$ & $5.391\,247(60) \times 10^{-44}$ s \\
Planck mass & $m_P$ & $2.176\,434(24) \times 10^{-8}$ kg \\
\textbf{Planck temperature} & $T_P$ & $\mathbf{1.416\,784(16)} \times 10^{32}$ K \cite{nist_tp} \\
Planck charge & $q_P$ & $1.875\,545\,956(41) \times 10^{-18}$ C \\
Fine-structure inverse & $\alpha^{-1}$ & $137.035\,999\,177(21)$ \cite{nist_alpha} \\
Boltzmann constant & $k_B$ & $1.380649 \times 10^{-23}$ J/K (exact) \\
Elementary charge & $e$ & $1.602176634 \times 10^{-19}$ C (exact) \\
Avogadro constant & $N_A$ & $6.02214076 \times 10^{23}$ mol$^{-1}$ (exact) \\
\bottomrule
\end{tabular}
\end{table}

Note: Four constants ($h$, $k_B$, $e$, $N_A$) are \textit{exact} by definition in the 2019 SI redefinition. Their digit sequences are therefore fixed, not subject to measurement uncertainty.

\subsection{Unified Digit Extraction}

For reproducibility, we define a unified digit extraction procedure:
\begin{enumerate}
    \item Take the CODATA 2022 recommended value mantissa (scientific notation).
    \item Extract digits up to the last significant digit before uncertainty dominates.
    \item For exact constants, extract all specified digits.
    \item Input sequence: decimal digits without decimal point, e.g., $T_P = 1.416\,784 \times 10^{32} \to [1,4,1,6,7,8,4]$.
\end{enumerate}

\textbf{Important}: Digit analysis is inherently representation-dependent (decimal, binary, hexadecimal) and unit-dependent. Results from digit analysis should be interpreted as \textit{computational curiosities}, not physical laws.

%===============================================================================
\section{Results}
%===============================================================================

\subsection{Summary of Findings}

\begin{table}[h]
\centering
\caption{Exploration results summary (corrected)}
\label{tab:results}
\begin{tabular}{@{}lllll@{}}
\toprule
Constant & Pattern & DoF & OOS & Classification \\
\midrule
$h$ & 2nd-order recurrence & 2 & $\times$ & DescriptiveModel \\
$\ell_P$ & 2nd-order recurrence & 2 & $\times$ & DescriptiveModel \\
$t_P$ & Linear recurrence & 1 & $\times$ & DescriptiveModel \\
$m_P$ & Linear recurrence & 1 & $\times$ & DescriptiveModel \\
$\mathbf{T_P}$ & \textbf{Quadratic} & \textbf{0} & $\checkmark$ & \textbf{PredictiveFit} \\
$q_P$ & 2nd-order recurrence & 2 & $\times$ & DescriptiveModel \\
$\mathbf{\alpha^{-1}}$ & \textbf{None} & \textbf{--} & $\times$ & \textbf{MYSTERY} \\
$k_B$ & Linear recurrence & 1 & $\times$ & DescriptiveModel \\
$e$ & 2nd-order recurrence & 2 & $\times$ & DescriptiveModel \\
$N_A$ & 2nd-order recurrence & 2 & $\times$ & DescriptiveModel \\
\bottomrule
\end{tabular}
\end{table}

\textbf{Key result}: No constant achieved \textbf{EmergentLaw} status. The most promising result ($T_P$) has DoF $= 0$, meaning OOS success \textbf{may be coincidental}.

\subsection{Planck Temperature: PredictiveFit (Not Discovery)}

The digit sequence of Planck temperature $T_P = 1.416\,784(16) \times 10^{32}$ K (CODATA 2022 \cite{nist_tp}):
\begin{equation}
    S_{T_P} = \{1, 4, 1, 6, 7, 8, 4\}
\end{equation}

\textbf{Complete Exploration Output} (verbatim from S-wind):

\begin{verbatim}
SELF-LOOP EXPLORATION:
  Iterations: 5, Patterns tried: 5
  Signal: 7 → 3 points
  Extensions: 0, Transformations: 4
  DoF improvement: +0
  ⏹️ EXHAUSTED: All patterns tried, no new
     signal generated
\end{verbatim}

\textbf{Exhaustion proof}: ``All patterns tried, no new signal generated'' means \textbf{closure reached}---every transformation leads to already-explored signals.

After 5 exploration iterations with 4 difference transformations, on a \textbf{3-point transformed signal}, S-wind found:

\begin{equation}
\boxed{f(n) = \frac{7}{18}n^2 - \frac{7}{6}n + \frac{26}{15}}
\end{equation}

\textbf{DoF Analysis}:
\begin{align}
    |S_{\text{fit}}| &= 3 \text{ (transformed signal points)} \\
    \dim(\Theta_f) &= 3 \text{ (parameters: } a, b, c \text{)} \\
    \text{DoF} &= 3 - 3 = 0
\end{align}

\textbf{Why DoF $= 0$ is problematic}:

By the polynomial interpolation theorem:
\begin{quote}
\textit{Any $n$ distinct points can be fitted exactly by a unique polynomial of degree $\leq n-1$.}
\end{quote}

Therefore, 3 points with a quadratic (3 parameters) is a \textbf{guaranteed exact fit}---it carries \textbf{no information about the underlying structure}.

\textbf{Verification}:
\begin{itemize}
    \item In-sample: 3/3 points (trivially guaranteed)
    \item Out-of-sample: 4/4 points matched ($n = 4, 5, 6, 7$)
    \item \textbf{Classification}: PredictiveFit (NOT Discovery)
\end{itemize}

\subsubsection{Why OOS Success May Be Coincidence}

The OOS verification passed, but we classify this as \textbf{PredictiveFit} rather than \textbf{EmergentLaw} because:

\begin{enumerate}
    \item \textbf{Zero constraint redundancy}: The original fit has DoF $= 0$. The quadratic formula was not ``constrained'' by the data---it was simply interpolated.
    
    \item \textbf{Insufficient sample size}: Only 7 digits total. Short digit sequences have high probability of exhibiting spurious patterns purely by chance.
    
    \item \textbf{No physical theory}: There is no theoretical reason why Planck temperature digits (in decimal, in SI units) should follow a quadratic law.
    
    \item \textbf{Transform-dependent}: The ``structure'' appears only after 4 difference transformations, not in the raw digits.
    
    \item \textbf{Statistical argument}: Given $10^7$ possible 7-digit sequences and multiple transformation chains, the probability of \textit{some} sequence showing OOS-passing quadratic structure is non-negligible.
\end{enumerate}

\textbf{Conclusion}: We cannot distinguish between:
\begin{itemize}
    \item (A) Genuine structure in $T_P$ digits
    \item (B) Statistical coincidence in a small sample
\end{itemize}

Conservative classification: \textbf{PredictiveFit} (testable hypothesis, not confirmed discovery).

\subsection{Fine-Structure Constant: MYSTERY (with Exhaustion Proof)}

The digit sequence of $\alpha^{-1} = 137.035\,999\,177(21)$ (CODATA 2022 \cite{nist_alpha}):
\begin{equation}
    S_{\alpha^{-1}} = \{1, 3, 7, 0, 3, 5, 9, 9, 9, 1, 7, 7\}
\end{equation}

\textbf{Complete Exploration Output} (verbatim from S-wind):

\begin{verbatim}
SELF-LOOP EXPLORATION:
  Iterations: 12, Patterns tried: 10
  Signal: 12 → 3 points
  Extensions: 0, Transformations: 11
  ✅ FATE CONVERGED: Final answer reached
\end{verbatim}

\subsubsection{Exhaustive Hypothesis Space (Auto-Generated)}

\begin{verbatim}
H = {
  H_algebraic.recurrence: s[n+1] = a·s[n] + b,
  H_algebraic.closed_form: f(n) = an + b, c/n, a·rⁿ,
  H_algebraic.identity: s[n]·n = k,
  H_algebraic.quadratic: f(n) = an² + bn + c,
  H_algebraic.cubic: f(n) = an³ + bn² + cn + d,
  H_algebraic.rational: f(n) = P(n)/Q(n),
  H_algebraic.multiplicative: f(mn) = f(m)·f(n),
  H_pattern.{constant, linear, quadratic, cubic,
             reciprocal, exponential, periodic,
             factorial, recurrence_2, recurrence_3,
             general_polynomial},
  H_oeis: 130 known sequences
}
\end{verbatim}

\textbf{Total}: $|H| = 19$ hypothesis classes + 130 OEIS sequences.

\subsubsection{Exhaustive Transformation Set (Auto-Generated)}

\begin{verbatim}
T = {
  T.identity: 12 signals,
  T.difference (Δ): 11 times,
  T.ratio: f(n+1)/f(n),
  T.cumsum: Σf(k),
  T.Δ: First difference (priority=100),
  T.|·|: Absolute value (priority=90),
  T.sign: Sign extraction (priority=85),
  T.Σ: Cumulative sum (priority=80),
  T.n·f: Scale by index (priority=70),
  T.√|·|: Square root of absolute (priority=65),
  T.MA3: Moving average (window 3) (priority=60),
  T.f²: Square transform (priority=55)
}
\end{verbatim}

\textbf{Total}: $|T| = 12$ transformation types applied iteratively until closure.

\subsubsection{Closure Reached = No Data Selection}

The critical proof that we did \textbf{not} select data:
\begin{equation}
\boxed{\text{CLOSURE}(E_S) := \forall \tau \in T, \forall S_i \in E_S: \text{sig}(\tau(S_i)) \in E_S}
\end{equation}

S-wind output confirms: ``\texttt{FATE CONVERGED: Final answer reached}'' means \textbf{every transformation} of \textbf{every explored signal} led to an already-explored signal. This is mathematical closure.

\textbf{Final state}:
\begin{itemize}
    \item 12 unique signals explored (closure set $|E_S| = 12$)
    \item Signal reduced to 3 points after full difference chain
    \item Best fit: Linear recurrence with DoF $= 1$
    \item OOS: FAILED at $n = 4$ (prediction failed)
\end{itemize}

\textbf{Classification}: \textbf{MYSTERY} (DoF $=$ --)

\subsubsection{Structure-Free Certificate (Bounded)}

\begin{center}
\fbox{\parbox{0.9\columnwidth}{
\textbf{Structure-Free Certificate}

\textbf{Sequence}: $S = \{1, 3, 7, 0, 3, 5, 9, 9, 9, 1, 7, 7\}$ (digits of $\alpha^{-1}$)

\textbf{Closure proof}: $|E_S| = 12$ signals, $|T| = 12$ transforms, $\text{CLOSURE}(E_S) = \text{TRUE}$

\textbf{Search bounds} (for certificate validity):
\begin{itemize}
    \item Hypothesis space: $|H| = 19$ classes + 130 OEIS
    \item Max degree: $d_{\max} = 11$
    \item Max recurrence order: $k_{\max} = 3$
    \item Transform depth: Full chain until length $< 3$
\end{itemize}

\textbf{Result}:
\begin{equation}
    \forall f \in H, \forall \tau \in T^*: \text{DoF}(f, \tau(S)) \leq 1 \text{ and OOS fails}
\end{equation}

\textbf{Conclusion}: No non-trivial algebraic structure found within $(H, T)$ under bounded closure exhaustion.
}}
\end{center}

\textbf{Note}: This is a \textit{bounded certificate}, not a mathematical theorem. It asserts exhaustion over a finite, explicitly defined and enumerated hypothesis space. Larger $(H', T')$ might find structure we missed---but within $(H, T)$, exhaustion is provable.

%===============================================================================
\section{Discussion}
%===============================================================================

\subsection{Why ``Theorem'' is Inappropriate}

We originally used ``Structure-Free Theorem'' but revised to ``Structure-Free Certificate'' because:

\begin{enumerate}
    \item A \textbf{theorem} implies a proof of an intrinsic property
    \item Our result is an \textbf{exhaustion statement} over $(H, T)$
    \item Larger $H$ or $T$ might find structure we missed
\end{enumerate}

The certificate is valuable because it:
\begin{itemize}
    \item Documents exactly what was searched
    \item Provides a reproducible negative result
    \item Defines the boundary of current exploration
\end{itemize}

\subsection{Metrological Considerations}

\subsubsection{Exact vs. Uncertain Constants}

Post-2019 SI redefinition:
\begin{itemize}
    \item $h$, $k_B$, $e$, $N_A$: \textbf{exact} (defining constants)
    \item $\ell_P$, $t_P$, $m_P$, $T_P$: derived, with uncertainty
    \item $\alpha^{-1}$: measured, with uncertainty $\pm 21 \times 10^{-9}$
\end{itemize}

For exact constants, digit analysis is well-defined. For measured constants, only digits within the uncertainty bound are meaningful.

\subsubsection{Unit Dependence}

Digit sequences depend on:
\begin{itemize}
    \item Choice of units (SI vs. CGS vs. natural units)
    \item Base (decimal vs. binary vs. hexadecimal)
\end{itemize}

Any ``structure'' found is contingent on these choices and may not reflect physical reality.

\subsection{Methodological Contributions}

\subsubsection{True Exhaustion via Closure Detection}

Unlike heuristic pattern-matching systems that use fixed iteration limits, S-wind guarantees \textbf{true exhaustion} through closure detection:

\begin{enumerate}
    \item \textbf{Signal signature}: Each explored signal is assigned a unique hash
    \item \textbf{Explored set}: $E = \{\text{sig}(S_0), \text{sig}(S_1), \ldots\}$
    \item \textbf{Closure condition}: $\forall \tau \in T: \text{sig}(\tau(S_i)) \in E$
\end{enumerate}

When no transform can produce a signal outside the explored set, the system has reached \textbf{closure}---all reachable signals have been tested. This transforms ``we didn't find anything'' into ``we exhaustively searched $(H, T)$ and found nothing.''

Typical exploration statistics:
\begin{itemize}
    \item 12-digit sequence: 12+ iterations, 11+ transforms, 12+ unique signals
    \item 15-digit sequence: 15+ iterations, 14+ transforms, 15+ unique signals
\end{itemize}

\subsubsection{Dynamic Container Generation}

S-wind's Space component analyzes signal characteristics:
\begin{itemize}
    \item Oscillation $\to$ try $|\cdot|$, sign extraction
    \item Non-monotonic $\to$ try cumulative sum
    \item Polynomial growth $\to$ try $\Delta^2$
\end{itemize}

This generates \textbf{signal-specific} transformation chains rather than fixed pipelines.

\subsubsection{Honest Classification}

The distinction between:
\begin{itemize}
    \item \textbf{EmergentLaw}: DoF $\geq 2$ + OOS pass (genuine constraint)
    \item \textbf{PredictiveFit}: DoF $= 0$ + OOS pass (may be coincidence)
\end{itemize}

prevents over-claiming discoveries from trivial interpolations.

%===============================================================================
\section{Algebraic Structure Analysis: Definition Derivation}
%===============================================================================

Beyond digit sequence analysis, S-wind's symbolic algebra system can \textbf{derive} the \textbf{intrinsic algebraic relations} between physical constants from their definitions. These relations are:
\begin{itemize}
    \item Independent of unit systems (SI, CGS, natural units)
    \item Independent of number bases (decimal, binary, hexadecimal)
    \item Independent of measurement precision
    \item Exact by definition, not empirical discoveries
\end{itemize}

\textbf{Key distinction}: Digit analysis seeks \textit{patterns in representations}; algebraic analysis derives \textit{relations from definitions}. The latter is more robust and physically meaningful.

\subsection{Algebraic Definitions (Input)}

By definition, the Planck units are expressed in terms of fundamental constants:
\begin{align}
    \ell_P &= \sqrt{\frac{\hbar G}{c^3}} & \text{(Planck length)} \\
    t_P &= \sqrt{\frac{\hbar G}{c^5}} & \text{(Planck time)} \\
    m_P &= \sqrt{\frac{\hbar c}{G}} & \text{(Planck mass)} \\
    T_P &= \sqrt{\frac{\hbar c^5}{G k_B^2}} & \text{(Planck temperature)} \\
    q_P &= \sqrt{4\pi\varepsilon_0 \hbar c} & \text{(Planck charge)} \\
    E_P &= m_P c^2 & \text{(Planck energy)}
\end{align}

The fine-structure constant:
\begin{equation}
    \alpha = \frac{e^2}{4\pi\varepsilon_0 \hbar c} = \frac{e^2}{q_P^2}
\end{equation}

\subsection{Derived Algebraic Relations (By Definition of $q_P$)}

S-wind's symbolic system derives the following \textbf{exact} algebraic relations directly from the definitions above:

\begin{enumerate}
    \item \textbf{Speed of light from Planck units}:
    \begin{equation}
        \boxed{\frac{\ell_P}{t_P} = c}
    \end{equation}
    
    \item \textbf{Action quantum}:
    \begin{equation}
        \boxed{\ell_P \cdot m_P \cdot c = \hbar}
    \end{equation}
    
    \item \textbf{Temperature-energy equivalence}:
    \begin{equation}
        \boxed{T_P \cdot k_B = m_P c^2 = E_P}
    \end{equation}
    
    \item \textbf{Fine-structure as charge ratio}:
    \begin{equation}
        \boxed{\alpha = \left(\frac{e}{q_P}\right)^2}
    \end{equation}
\end{enumerate}

\subsection{Dimensionless Combinations}

The following dimensionless combinations are exactly unity:
\begin{align}
    \frac{\ell_P \cdot m_P \cdot c}{\hbar} &= 1 \\
    \frac{t_P \cdot E_P}{\hbar} &= 1 \\
    \frac{T_P \cdot k_B}{E_P} &= 1 \\
    \alpha \cdot \frac{q_P^2}{e^2} &= 1
\end{align}

\subsection{Digit Analysis vs. Algebraic Analysis}

\begin{table}[h]
\centering
\caption{Comparison of analysis approaches}
\begin{tabular}{@{}lll@{}}
\toprule
Aspect & Digit Analysis & Algebraic Analysis \\
\midrule
Input & $\{1,4,1,6,7,8,4\}$ & $T_P = \sqrt{\hbar c^5/G k_B^2}$ \\
Dependence & Base, units, precision & \textbf{None} \\
Output & Statistical fit & \textbf{Exact relations} \\
Physical meaning & Weak (possible coincidence) & \textbf{Strong} \\
Discovery type & PredictiveFit (DoF=0) & \textbf{Verified identity} \\
\bottomrule
\end{tabular}
\end{table}

\textbf{Key insight}: While digit analysis of $T_P$ yields only a PredictiveFit with DoF $= 0$, algebraic analysis reveals that $T_P$ participates in the \textbf{exact} relation $T_P \cdot k_B = E_P$. This is a genuine physical structure, not a statistical artifact.

%===============================================================================
\section{Conclusion}
%===============================================================================

We have presented a systematic analysis of 10 fundamental physical constants using the S-wind framework with CODATA 2022 values, employing both digit sequence analysis and symbolic algebraic analysis.

\textbf{Main findings}:

\begin{enumerate}
    \item \textbf{Digit analysis}: No EmergentLaw discovered. All constants either have DoF $= 0$ or fail OOS verification. The $T_P$ quadratic fit (DoF $= 0$, OOS pass) is classified as \textbf{PredictiveFit}---possible coincidence.
    
    \item \textbf{Algebraic analysis}: S-wind's symbolic system verifies \textbf{exact} algebraic relations between Planck units:
    \begin{itemize}
        \item $\ell_P \cdot m_P \cdot c = \hbar$ (action quantum)
        \item $T_P \cdot k_B = E_P$ (temperature-energy equivalence)
        \item $\alpha = (e/q_P)^2$ (fine-structure as charge ratio)
    \end{itemize}
    These are \textbf{genuine physical structures}, independent of units or number base.
    
    \item \textbf{Structure-Free Certificates}: For digit sequences, we introduce formal certificates with explicit hypothesis space $H$, transformation set $T$, and closure guarantee.
    
    \item \textbf{Methodological contrast}: Digit analysis yields statistical fits (weak); algebraic analysis yields exact identities (strong). Future work should prioritize algebraic exploration.
\end{enumerate}

\textbf{Recommendation}: For physical constants, \textbf{algebraic structure analysis} (exploring relations like $\ell_P m_P c = \hbar$) is more meaningful than digit sequence analysis. S-wind supports both approaches.

The S-wind framework and all experimental data are available at \url{https://github.com/rogueintelligence/s-wind}.

%===============================================================================
\begin{acknowledgments}
This work was performed by the S-wind automated discovery system with human oversight for methodological rigor. Physical constant values sourced from NIST CODATA 2022.
\end{acknowledgments}

\begin{thebibliography}{9}
\bibitem{planck1900} M. Planck, Ann. Phys. \textbf{4}, 553 (1901).
\bibitem{dirac1937} P. A. M. Dirac, Nature \textbf{139}, 323 (1937).
\bibitem{barrow2002} J. D. Barrow, \textit{The Constants of Nature} (Pantheon, 2002).
\bibitem{bailey1997} D. H. Bailey et al., Math. Comp. \textbf{66}, 903 (1997).
\bibitem{tegmark2014} M. Tegmark, \textit{Our Mathematical Universe} (Knopf, 2014).
\bibitem{codata2022} P. J. Mohr, E. Tiesinga, D. B. Newell, and B. N. Taylor, Rev. Mod. Phys. \textbf{97}, 025002 (2025). [CODATA 2022 recommended values]
\bibitem{nist_constants} NIST, ``Fundamental Physical Constants,'' \url{https://physics.nist.gov/cuu/Constants/} (accessed 2026-01-25).
\bibitem{nist_tp} NIST, ``Planck temperature,'' \url{https://physics.nist.gov/cgi-bin/cuu/Value?plktmp} (accessed 2026-01-25).
\bibitem{nist_alpha} NIST, ``Inverse fine-structure constant,'' \url{https://physics.nist.gov/cgi-bin/cuu/Value?alphinv} (accessed 2026-01-25).
\end{thebibliography}

\end{document}
